%================================================================

\subsection{Definition}

Define $p(x)\in\Qx$ to be prime if:

\begin{itemize}
    \item $\forall a(x),b(x)\in\Qx$ where $a(x)b(x)=p(x)$, either $a(x)$ or $b(x)$ is a unit
    \item $p(x)$ is not a unit
    \item where $p(x)=\sum_{0}^{n} a_i x^i$, $a_n=1$ (p(x) is monic)
\end{itemize}

Note that all linears $l(x)\in\Qx$ (polynomials of degree 1) are prime in $\Qx$ since they have a norm of 2. for any $a(x),b(x)\in\Qx$ where $a(x)b(x)=l(x)$, $\|a(x)\|\|b(x)\|=2$ (norm preserves multiplication), and $\|a(x)\|,\|b(x)\|\in\bZ^+$ (since they cannot be 0). Then, because 2 is prime in $\bZ^+$, $\|a(x)\|=1\lor\|b(x)\|=1$. This means that either $a$ or $b$ is a unit, and it follows then that $l(x)$ is a prime.

%================================================================

\subsection{Prime Factors}

\begin{theorem}[Existence of prime factors]
    $\forall a(x)\in\Qx$ where $a(x)$ isn't a unit, $\exists p(x)$ which is prime where $p(x)|a(x)$.\\
\end{theorem}
    

When $a(x)=0$, it is divisible by all primes. For some polynomial $a(x)$ which isn't 0 or a unit, let $S=\{n(x)\in\Qx\mid n(x)\neq0\land\|{n(x)}\|\neq1\land n(x)\mid a(x)\}$. We know S is non-empty since $a(x)\in S$.\\

Then consider $\|S\|=\{\|n(x)\|\in S\}$. Since $\forall n(x)\in\Qx$ where $n\neq0$, $\|n(x)\|\in\bZ^+$, $\|S\|\in\bZ^+$. Then, since it is non-empty, there must exist an element $\|m(x)\|\in\|S\|$ which corresponds to $m(x)\in S$ where $\|m(x)\|$ is minimal in $\|S\|$.\\

Suppose $\exists k(x)\mid m(x)$, where $\|k(x)\|\neq\|m(x)\|$ and isn't a unit. Since $\|k(x)\|\mid \|m(x)\|$ and $\|k(x)\|,\|m(x)\|\in\bZ^+$, $ \|k(x)\|\leq \|m(x)\|$, but $\|k(x)\|\neq \|m(x)\|$, so $\|k(x)\|<\|m(x)\|$. Since $k(x)\mid m(x)$, then because $m(x)\mid a(x)$, $k(x)\mid a(x)$, so $k(x)\in|S$. Then, we have $\|k(x)\|\in\|S\|$ too, but $\|k(x)\|<\|m(x)\|$, so $\|m(x)\|$ is not minimal in $S\xrightarrow{}\xleftarrow{}$.\\

Therefore, if $\exists k(x)\mid m(x)$, $\|k(x)\|=\|m(x)\|$ or it is a unit. For any $a(x)b(x)=m(x)$, either $a(x)$ is a unit or $\|a(x)\|=\|m(x)\|$. In the latter case, this implies that $b(x)$ is a unit, since preservation of norm after multiplication is unique to units. Consequently, $\forall a(x),b(x)\in\Qx$ where $a(x)b(x)=m(x)$, either $a(x)$ or $b(x)$ is a unit, and $m(x)$ isn't a unit, meaning that $m(x)$ is prime.\\

%================================================================

\subsection{Prime Factorisation}

\begin{theorem}[Existence of a Prime Factorisation]
    $\forall n(x) \in \Qx$, $n(x)$ can be expressed as $u(x)\prod _{i=1}^n p_i(x)^{\alpha_i}$ where $p_i(x) \in \Qx$ is prime, $a_i\in\bZ^+$, and $u(x) \in \Qx^{\times}$.
\end{theorem}

Let $S= \{(p_i, \alpha_i)\}$. $\forall n(x)\in\bZ^+,\exists m_1(x)\mid n(x)$ and $m_1(x)$ is prime. If $m_1(x)=n(x)$, then such $S$ exists where $S=\{m_1\}$. Otherwise, $\exists n_1(x)$ where $m_1(x)n_1(x)=n(x)$. It follows then that $\|m_1(x)\|\|n_1(x)\|=\|n(x)\|$, so $\|n_1(x)\|\mid\|n(x)\|$. Since they are both positive integers, this means that $\|n_1(x)\|\leq\|n(x)\|$. However, if $\|n_1(x)\|=\|n(x)\|$, then $\|m_1(x)\|=1$, meaning it is a unit and therefore can't be prime. Therefore, $\|n_1(x)\|<\|n(x)\|$.\\

Similarly, either $n_1(x)$ is prime or $m_2(x)n_2(x)=n_1(x)$, where $m_2(x)$ is prime. If this terminates at $n_k(x)$ which is a unit, we have $n(x)=n_k(x)\prod_{i=1}^{k}m_i(x)$. Then, we can have $S'=\{\{m_i(x)\mid i\in\bZ \land 0<i\leq k\}\}$. Where $c(m_i(x))$ returns its multiplicity within $S'$, we can have $S=\{(m_i(x),c(m_i(x))\mid m_i(x)\in S'\}$, which matches the properties required in our proposition.\\

Suppose this doesn't terminate. Then, we have set $\|N\|=\{\|n_i(x)\||i\in\bZ^+\}$ which is infinite, and $\forall k\in\bZ^+, \|n_k(x)\|>\|n_{k+1}(x)\|$. Also, $\|N\|\subset\bZ^+$ since the norm of all non-zero polynomials is a positive integer. Since it is non-empty, there must be a minimal element $\|n_m(x)\|$ in it. However, $\|n_m(x)\|>\|n_{m+1}(x)\|$, so it isn't minimal $\xrightarrow{}\xleftarrow{}$.\\
 
Therefore, it must terminate, in which case there exists such set $S$.\\
 
%================================================================
